
\section{Persistent parameters and particle approximation}
A particle approximation can reduce computation after a probability model is
chosen. It does not create that model. A type with probability $0.01$ has
probability $0.99^{100}>1/3$ of disappearing from 100 independent multinomial
resampling draws. No new evidence is involved in this loss. Rare but material
types therefore need explicit support checks, stratification or separate
conditional computations, with the resulting approximation stated.

For state-space models with persistent unknown parameters, SMC$^2$ combines
parameter particles and state filters, with particle-MCMC rejuvenation
\cite[Section~3.2, equation~(7), Proposition~1, Appendix~A]{litclosesmc2}. Its
extended target has the desired parameter posterior as a marginal under the
source assumptions. Finite particles still approximate that posterior; exactness
of a target marginal does not mean that one finite run represents every
possible type. The unbiased full particle likelihood estimate used in this
construction must be distinguished from an arbitrary estimated likelihood
ratio.

The inspected author-linked modern implementation creates a state filter for
each parameter particle, runs that filter through the current observation date
when evaluating a proposed parameter, and incorporates a likelihood correction
when increasing the state-particle count. These operations have real time and
memory costs. Retaining only filtering arrays differs from retaining complete
genealogies. This is a candidate for explicitly probabilistic filtering, not
a replacement for symbolic unweighted historical ambiguity.
